\documentclass{article}

\usepackage[margin=0.5in]{geometry}
\usepackage{amsmath, mathtools, array, booktabs, hyperref, amssymb}

\title{HW 4}

\author{Darin Peries}

\date{\today}

\begin{document}
\maketitle

\href{URL}{text}

\section*{Problem 1}

Prove that $\|v\|_B := \sqrt{v^*Bv}$ defines a norm on $\mathbb{C}^n$ given that $B$ is both positive definite and hermitian. \\

\subsection*{Proof for homogeneity}

Given $\lambda \in \mathbb{C}$ and $x \in \mathbb{C}^n$. $\|\lambda x\|_B = \sqrt{(\lambda^* x^*)B(\lambda x)} = \sqrt{|\lambda|^2 x^*Bx}=|\lambda|\sqrt{x^*Bx}=|\lambda|\|x\|_B$.
Therefore the norm maintains homogeneity.

\subsection*{Proof for triangle inequality}

Given $x,\ y\in\mathbb{C}^n$, $\|x + y\|_B = \sqrt{(x+y)^*B(x+y)} = \sqrt{(x^* + y^*)B(x + y)} = \sqrt{(x^*B + y^*B)(x+y)} = \sqrt{x^*Bx + x^*By + y^*Bx + y^*By}$. Because
matrix $B$ is hermitian, $(x^*By)^* = y^*Bx$. Therefore, the expression can be rewritten as $\sqrt{x^*Bx + 2\operatorname{Re}(y^*Bx) + y^*By}$. Using the fact that $\operatorname{Re}(z) \leq |z|$ and the Cauchy-Schwarz inequality 
for the $B$-inner product (which works since conjugate symmetry holds due to hermitian and positive definiteness gives positivity), $|\langle a, \ b \rangle_B| \leq \|a\|_B \|b\|_B = \sqrt{a^*Ba} \sqrt{b^*Bb}$, given that $|\langle a, \ b \rangle_B| = |a^*Bb|$. The expression can further be rewritten
as 
\[ \begin{gathered} \sqrt{x^*Bx + 2\operatorname{Re}(y^*Bx) + y^*By} \leq \sqrt{x^*Bx + 2|y^*Bx| + y^*By} \leq \\
 \sqrt{x^*Bx + 2\sqrt{y^*By}\sqrt{x^*Bx} + y^*By} = \sqrt{(\sqrt{x^*Bx} + \sqrt{y^*By})^2} = \sqrt{(\|x\|_B + \|y\|_B)^2}=\|x\|_B + \|y\|_B \end{gathered} \]

\noindent From the expression above, $\|x+y\|_B \leq \|x\|_B + \|y\|_B$ and the triangle inequality hold.

\subsection*{Proof for Non-negativity and definiteness}

Since $B$ is hermitian, $(x^*Bx)^* = x^*B^*x^{**}= x^*Bx$, and becuase $B$ is positive definite, $x^*Bx$ is strictly greater than zero and real.
for $x\neq 0$. Therefore, $\|x\|_B$ is positive and real for $x \neq 0$ and $\|0\|_B=0$. Non-negativity and definiteness is upheld.\\

\noindent From the three proofs above, $\| \cdot \|_B$ is a proper norm in $\mathbb{C}^n$.

\section*{Problem 2}

Given an A-orthonormal system, prove that $A$ is symmetric-positive definite.\\ 

\noindent Defining a matrix $U\in\mathbb{R}^{n\times n}$ using the column vectors of the A-orthonormal system, $U^TAU=I$, 
as the multiplication of the A-orthonormal matrices and $A$ should produce the identity matrix. Additionally, If $\sum_i c_iu(i) = 0$, multiplying on the left by $u(j)^TA$ gives $c_j = 0$ for each $j$.
So the columns of $U$ are independent, and $U$ is invertible. By transposing $$(U^TAU)=I=I^T=(U^TAU)^T=U^TA^TU=U^TAU$$ 
Because the equation holds, and the fact that $U$ is invertible due to being a square matrix and having linearly independent column vectors, 
multiplying the prior equation by $U^{-T}$ on the left and $U^{-1}$ on the right gives $A=A^T$. Therefore, $A$ is symmetric. 
Say for all $x\in \mathbb{R}^n$ and $x \neq 0$, $v=U^{-1}x; \ x = Uv$. $x^TAx=v^TU^TAUv=v^TIv=v^Tv>0$. Therefore $A$ is positive definite.


\section*{Problem 3}

Prove that $2Q_A(v) = (e,e)_A - (u,u)_A$, given that $Q_A(v) := \frac{1}{2}(v,v)_A - (b,v)_I$ and $Au=b, \ (b,v)_I=b^Tv$\\

\noindent $2Q_A(v) = (v,v)_A - 2(b,v)_I = (v,v)_A - 2(Au, v)_I$. Because $(b,v)_I = b^Tv$, $(Au, v)_I = (Au)^Tv = u^TA^Tv = u^TAv = (u, v)_A$ due to symmetry.
Therefore, the expression can be written as $(v,v)_A - 2(u,v)_A$. By completing the square through adding and subtracting $(u,u)_A$, the new 
expression is 
$$ (v,v)_A - 2(u,v)_A + (u,u)_A - (u,u)_A = (v-u,v-u)_A - (u,u)_A = (e,e)_A - (u,u)_A $$.

\section*{Problem 4}

\subsection*{Part 1}

\begin{table}[ht]
\centering
\caption{Approximate solution vectors obtained with conjugate gradient (CG)
and gradient descent (GD) for both matrices $A_1, \ A_2$.}
\label{tab:solution-vectors}
\renewcommand{\arraystretch}{1.3}
\begin{tabular}{c c p{11cm}}
\hline
Method & Matrix & Solution vector $x^T$ \\
\hline

CG & $A_1$ &
\texttt{[0.333, 0.333, 0.333, 0.333, 0.333, 0.333, 0.333, 0.333,} \newline
\texttt{ 0.333, 0.333, 0.333, 0.333, 0.333, 0.333, 0.333, 0.333]} \\[0.75em]

CG & $A_2$ &
\texttt{[0.333, 0.333, 0.333, 0.333, 0.333, 0.333, 0.333, 0.333,} \newline
\texttt{ 0.333, 0.333, 0.333, 0.333, 0.333, 0.333, 0.333, 0.333,} \newline
\texttt{ 0.333, 0.333, 0.333, 0.333, 0.333, 0.333, 0.333, 0.333,} \newline
\texttt{ 0.333, 0.333, 0.333, 0.333, 0.333, 0.333, 0.333, 0.333]} \\[0.75em]

GD & $A_1$ &
\texttt{[0.333, 0.333, 0.335, 0.333, 0.332, 0.332, 0.332, 0.334,} \newline
\texttt{ 0.335, 0.335, 0.335, 0.335, 0.334, 0.333, 0.332, 0.330]} \\[0.75em]

GD & $A_2$ &
\texttt{[0.333, 0.333, 0.332, 0.334, 0.335, 0.334, 0.334, 0.333,} \newline
\texttt{ 0.332, 0.332, 0.332, 0.332, 0.332, 0.333, 0.333, 0.334,} \newline
\texttt{ 0.334, 0.334, 0.335, 0.335, 0.335, 0.335, 0.335, 0.335,} \newline
\texttt{ 0.335, 0.334, 0.334, 0.333, 0.332, 0.332, 0.331, 0.330]} \\
\hline
\end{tabular}
\end{table}

\subsection*{Part 2}

\begin{table}[ht]
\centering
\caption{Solutions obtained using different iterative methods.}
\label{tab:iterative-solutions}
\begin{tabular}{c c}
\hline
Method & Approximate solution $x^T$ \\
\hline
Jacobi &
\texttt{[1.0, -2.0, 3.0, -2.0, 1.0]} \\

Gauss--Seidel &
\texttt{[1.0, -2.0, 3.0, -2.0, 1.0]} \\

Gradient Descent &
\texttt{[1.0, -2.0, 3.0, -2.0, 1.0]} \\

Conjugate Gradient &
\texttt{[1.0, -2.0, 3.0, -2.0, 1.0]} \\
\hline
\end{tabular}
\end{table}

\end{document}